\chapter{The Language Model Layer}

Four client classes live in \fpath{src/llm/}. They target different APIs but expose the same two methods, so any of them can be dropped into an agent without the agent knowing which one it got.

\section{The informal interface}

There is no abstract base class here --- the contract is duck-typed. Every client implements:

\begin{lstlisting}[style=py]
def generate(self, prompt: str,
             system_prompt: Optional[str] = None,
             temperature: float = ...,
             max_tokens: int = ...,
             conversation_history: Optional[List[Dict]] = None,
             **kwargs) -> str

def generate_with_context(self, query: str, context: str,
                          system_prompt: str,
                          temperature: float = ...,
                          max_tokens: int = ...,
                          conversation_history: Optional[List[Dict]] = None) -> str
\end{lstlisting}

\code{generate\_with\_context} is the RAG entry point and is the same in all four files --- it wraps the retrieved passages and the question in a fixed template, then delegates to \code{generate}:

\begin{lstlisting}[style=py]
        prompt = f"""Context from Ayurvedic texts:

{context}

---

User Query: {query}

Based on the context provided above, please provide a response."""
\end{lstlisting}

Specialist agents call \code{generate\_with\_context}. The orchestrator calls \code{generate} directly, both for the fast path and for synthesis, because by then the context has already been consumed and what it is passing along is the specialists' prose.

\section{The four clients side by side}

\begin{table}[H]
\centering
\small
\begin{tabularx}{\textwidth}{@{}l l l X@{}}
\toprule
\textbf{Client} & \textbf{Default model} & \textbf{Selected by} & \textbf{Notes} \\
\midrule
\code{GoogleClient} & \code{gemini-1.5-flash} & fallback, always last & The de facto default. \code{google-generativeai} SDK, stateless \code{generate\_content}. \\
\addlinespace
\code{OpenAIClient} & \code{gpt-4o-mini} & \code{USE\_OPENAI=true} & \code{openai} SDK. Thinnest of the four at 74 lines. \\
\addlinespace
\code{OllamaClient} & \code{llama3.2:3b} & \code{USE\_LOCAL\_LLM=true} & Raw HTTP to \code{localhost:11434}. Free and offline; ignores \code{max\_tokens}. \\
\addlinespace
\code{OpenRouterClient} & \code{mistral-7b-instruct:free} & \textbf{never} & Raw HTTP. Not imported anywhere in the application. \\
\bottomrule
\end{tabularx}
\caption{Client comparison. Only the first three are reachable from the running application.}
\end{table}

\section{Provider selection}

\code{AyurMindApp.\_\_init\_\_} walks a fixed preference order, catching failures at each step:

\begin{figure}[H]
\centering
\begin{tikzpicture}[node distance=7mm]
  \node[dia=clay, text width=22mm] (d1) {USE\_LOCAL\_LLM};
  \node[blk=midgreen, right=20mm of d1, minimum width=30mm] (ol) {OllamaClient};
  \node[dia=clay, below=13mm of d1, text width=22mm] (d2) {USE\_OPENAI};
  \node[blk=midgreen, right=20mm of d2, minimum width=30mm] (oa) {OpenAIClient};
  \node[blk=deepgreen, below=13mm of d2, minimum width=30mm] (gg) {GoogleClient};
  \node[blk=brick, below=9mm of gg, minimum width=30mm] (er) {raise};

  \draw[arb=midgreen] (d1) -- node[lbl,above]{true} (ol);
  \draw[arb=clay] (d1) -- node[lbl,right]{false} (d2);
  \draw[arb=midgreen] (d2) -- node[lbl,above]{true} (oa);
  \draw[arb=clay] (d2) -- node[lbl,right]{false} (gg);
  \draw[arb=brick] (gg) -- node[lbl,right]{no key} (er);

  \draw[ar, dashed] (ol.south) to[out=-90,in=0] node[lbl,right,yshift=1mm]{init failed} (d2.east);
  \draw[ar, dashed] (oa.south) to[out=-90,in=0] node[lbl,right,yshift=1mm]{init failed} (gg.east);
\end{tikzpicture}
\caption{Fallback chain. A client that raises during construction is logged as a warning and the next option is tried; only a Google failure is fatal.}
\end{figure}

\begin{lstlisting}[style=py,firstnumber=48]
        self.llm_client = None
        use_local_llm = os.getenv("USE_LOCAL_LLM", "false").lower() == "true"
        use_openai    = os.getenv("USE_OPENAI",    "false").lower() == "true"

        if use_local_llm:
            try:
                self.llm_client = OllamaClient()
                app_logger.info(f"Using Local client: {self.llm_client.model_name}")
            except Exception as e:
                app_logger.warning(f"Local LLM Client failed: {e}. Falling back.")

        if self.llm_client is None and use_openai:
            try:
                self.llm_client = OpenAIClient()
                app_logger.info(f"Using OpenAI client: {self.llm_client.model_name}")
            except Exception as e:
                app_logger.warning(f"OpenAI Client failed: {e}. Falling back to Google.")

        if self.llm_client is None:
            try:
                self.llm_client = GoogleClient()
                app_logger.info(f"Using Google client: {self.llm_client.model_name}")
            except Exception as e:
                app_logger.error(f"Failed to initialize GoogleClient: {e}.")
                raise e     # Critical failure if no client can be initialized
\end{lstlisting}

The fallback covers construction only. Once a client is chosen, a mid-session API error propagates up to the Gradio handler and is shown as an error message in the chat; there is no runtime failover to a second provider.

Note also that \code{OllamaClient.\_\_init\_\_} does not raise when Ollama is unreachable --- it logs a warning and returns a usable object. So \code{USE\_LOCAL\_LLM=true} with Ollama stopped selects the Ollama client successfully and then fails on the first generation call, rather than falling through to OpenAI or Google.

\section{Conversation history handling}

All four clients accept a \code{conversation\_history} list of \code{\{"role", "content"\}} dictionaries, and this is where their behaviour genuinely diverges.

\subsection{Google}

Gemini names the assistant role \code{model}, so history is remapped. The client also has to cope with Gradio's message format, where \code{content} can be a list of part-dictionaries rather than a plain string:

\begin{lstlisting}[style=py,firstnumber=62]
    def _format_history(self, conversation_history):
        if not conversation_history:
            return []
        formatted_history = []
        for turn in conversation_history:
            # Gemini uses 'model' for the assistant's role, 'user' for user
            role = "model" if turn["role"] == "assistant" else "user"
            text_content = self._extract_text_content(turn["content"])
            formatted_history.append({"role": role, "parts": [text_content]})
        return formatted_history
\end{lstlisting}

\code{\_extract\_text\_content} handles four shapes --- a plain string, a list of dicts with \code{text} keys, a single dict with a \code{text} key, and anything else via \code{str()}. That defensiveness exists because Gradio has changed its chat message schema more than once.

\subsection{OpenRouter and Ollama drop the last turn}

Both slice the history before using it:

\begin{lstlisting}[style=py]
        # OpenRouterClient.generate
        if conversation_history:
            # Add all but the last user message from history
            for turn in conversation_history[:-1]:
                messages.append(turn)
\end{lstlisting}

That slice is correct for how the UI calls them. \code{AyurMindApp.chat()} appends the user's new message to \code{history} \emph{before} passing it down, so the last element of \code{conversation\_history} is the same message that arrives separately as \code{prompt}.

\begin{bug}[title={Google and OpenAI send the current message twice}]
\code{GoogleClient.\_format\_history} iterates the whole list, then appends the prompt on top. \code{OpenAIClient} does the same. Since the UI has already appended the user's message to \code{history}, the model receives it once verbatim from history and again wrapped in the RAG or synthesis prompt. With the default Google client this happens on every slow-path turn --- once per specialist agent plus once for synthesis. Adding \code{[:-1]} in both clients, matching the OpenRouter and Ollama behaviour, resolves it.
\end{bug}

\subsection{Ollama flattens history into the prompt}

The Ollama \code{/api/generate} endpoint takes one string, so the client builds the transcript by hand:

\begin{lstlisting}[style=py,firstnumber=66]
        full_prompt = ""
        if system_prompt:
            full_prompt += f"{system_prompt}\n\n"
        if history_str:
            full_prompt += f"Previous conversation:\n{history_str}\n"
        full_prompt += f"Current query: {prompt}"

        payload = {
            "model": self.model,
            "prompt": full_prompt,
            "stream": False,
            "options": {"temperature": temperature,
                        "num_predict": 800}      # Ollama's max_tokens equivalent
        }
\end{lstlisting}

\code{num\_predict} is fixed at 800 regardless of what the caller asked for, and \code{max\_tokens} is popped from \code{kwargs} and discarded. Agents request 3000 tokens and the orchestrator requests 4000; against Ollama they all get 800. That truncation explains a great deal about the local-model evaluation outputs in Chapter~\ref{ch:eval} --- the four-section report the synthesis prompt asks for does not fit in 800 tokens.

\section{Error handling}

\code{OpenRouterClient} is the only one that translates HTTP status codes into actionable messages:

\begin{lstlisting}[style=py,firstnumber=70]
        except requests.exceptions.HTTPError as e:
            status_code = e.response.status_code
            if status_code == 401:
                raise RuntimeError("OpenRouter API Error: Invalid API Key. ...")
            elif status_code == 402:
                raise RuntimeError("OpenRouter API Error: You have exceeded "
                                   "your free tier limit or credits.")
            elif status_code == 404:
                raise RuntimeError(f"OpenRouter API Error: Model '{self.model}' "
                                   "not found. ...")
\end{lstlisting}

The other three wrap everything in a single \code{RuntimeError} carrying the original exception text. Google additionally treats an empty \code{response.text} as an error, which catches the case where Gemini's safety filters block a generation and return a response with no candidates --- a realistic outcome when the prompt describes symptoms and asks for treatment.

\begin{lstlisting}[style=py,firstnumber=142]
            if not response.text:
                raise ValueError("Empty response from model")
\end{lstlisting}

\section{Timeouts}

Both HTTP clients use 180 seconds. That is generous, and deliberately so: the slow path can issue four sequential generations, and a 3B model on CPU is not fast.

\begin{table}[H]
\centering
\small
\begin{tabular}{@{}l l l@{}}
\toprule
\textbf{Client} & \textbf{Timeout} & \textbf{Set where} \\
\midrule
\code{OpenRouterClient} & 180~s & \code{requests.post(..., timeout=180)} \\
\code{OllamaClient} & 180~s generate, 2~s availability probe & \code{is\_available()} uses the short one \\
\code{GoogleClient} & SDK default & not overridden \\
\code{OpenAIClient} & SDK default & not overridden \\
\bottomrule
\end{tabular}
\caption{Timeout settings. A full slow-path consultation can legitimately take several minutes.}
\end{table}
